% GENERATED FILE. Edit canonical lessons, then run npm run generate:books.
% canonical-source-hash: fnv1a64:c5cdab66593d3995
% canonical-lessons: PT-C305-ambulancia, PT-C305-sirene, PT-C305-socorrer, PT-C305-socorrista, PT-C305-maca

\chapter{Chegou a ambulância}
\label{ch:pt-chegou-a-ambulancia}
\hlchaptermodality{305}

\begin{chapteropening}
\textbf{By the end of this chapter:} \emph{I can describe help arriving: the ambulance and its siren, the first-aiders and the stretcher.}

Complete the last lesson of chapter 305: I can describe help arriving: the ambulance and its siren, the first-aiders and the stretcher.
\end{chapteropening}

\section[a ambulância]{a ambulância — the ambulance}

\label{lesson:PT-C305-ambulancia}



\begin{quote}
Before the new one: say the Portuguese for the landslide, then the Portuguese for to bury under earth or rubble.
\end{quote}

\subsection*{You'll want to know: a ambulância}
\textbf{a ambulância} — \textquotedblleft{}the ambulance\textquotedblright{}.

\begin{quote}
A ambulância chegou em dez minutos.
\end{quote}

\begin{quote}
\emph{The ambulance arrived in ten minutes.}
\end{quote}

\emph{Chamar uma ambulância} — to call an ambulance.

\begin{cousinweb}
\emph{ambulância} came from French \emph{ambulance} (EtymDB), first a field hospital that moved with an army, built on \emph{ambulant} \textquotedblleft{}walking\textquotedblright{}, from Latin \emph{ambulāre} \textquotedblleft{}to walk\textquotedblright{}.
\end{cousinweb}

\begin{grammarlens}[title={What you've built}]
One more word for telling, explaining and describing at B1.
\end{grammarlens}

\subsection*{Your turn}
\begin{itemize}
\raggedright
  \item \emph{Say it:} \emph{a ambulância}
  \item \emph{Say it:} \emph{a ambulância}, once more
  \item \emph{Say it:} a sentence of your own with \emph{a ambulância}
  \item \emph{Recall:} say the Portuguese for the landslide, then the Portuguese for to bury under earth or rubble, then say \emph{a ambulância} again
\end{itemize}

\begin{tcolorbox}[breakable,colback=teal!4,colframe=teal!35!black,title={Before you move on}]
What does a ambulância mean? (\textquotedblleft{}The ambulance\textquotedblright{}.) Say it once more.
\end{tcolorbox}
\section[a sirene]{a sirene — the siren}

\label{lesson:PT-C305-sirene}



\begin{quote}
Before the new one: say the Portuguese for to bury under earth or rubble, then the Portuguese for the ambulance.
\end{quote}

\subsection*{You'll want to know: a sirene}
\textbf{a sirene} — \textquotedblleft{}the siren\textquotedblright{}.

\begin{quote}
Ouvi a sirene de uma ambulância.
\end{quote}

\begin{quote}
\emph{I heard an ambulance siren.}
\end{quote}

Also \emph{a sirena}.

\begin{cousinweb}
\emph{sirene} was borrowed through French \emph{sirène} from Latin \emph{sīrēn}, from Greek \emph{seir\'{\={e}}n}, the singing creature of old stories (Lewis \& Short).
\end{cousinweb}

\begin{grammarlens}[title={What you've built}]
One more word for telling, explaining and describing at B1.
\end{grammarlens}

\subsection*{Your turn}
\begin{itemize}
\raggedright
  \item \emph{Say it:} \emph{a sirene}
  \item \emph{Say it:} \emph{a sirene}, once more
  \item \emph{Say it:} a sentence of your own with \emph{a sirene}
  \item \emph{Recall:} say the Portuguese for to bury under earth or rubble, then the Portuguese for the ambulance, then say \emph{a sirene} again
\end{itemize}

\begin{tcolorbox}[breakable,colback=teal!4,colframe=teal!35!black,title={Before you move on}]
What does a sirene mean? (\textquotedblleft{}The siren\textquotedblright{}.) Say it once more.
\end{tcolorbox}
\section[socorrer]{socorrer — to help (someone in danger), to come to the aid of}

\label{lesson:PT-C305-socorrer}



\begin{quote}
Before the new one: say the Portuguese for the ambulance, then the Portuguese for the siren.
\end{quote}

\subsection*{You'll want to know: socorrer}
\textbf{socorrer} — \textquotedblleft{}to help (someone in danger), to come to the aid of\textquotedblright{}.

\begin{quote}
Um rapaz socorreu a senhora que caiu.
\end{quote}

\begin{quote}
\emph{A young man helped the lady who fell.}
\end{quote}

\emph{Socorro!} is the cry for help.

\begin{cousinweb}
\emph{socorrer} is inherited from Latin \emph{succurrere} \textquotedblleft{}to run to help\textquotedblright{} (EtymDB), from \emph{sub} and \emph{currere} \textquotedblleft{}to run\textquotedblright{}.
\end{cousinweb}

\begin{grammarlens}[title={What you've built}]
One more word for telling, explaining and describing at B1.
\end{grammarlens}

\subsection*{Your turn}
\begin{itemize}
\raggedright
  \item \emph{Say it:} \emph{socorrer}
  \item \emph{Say it:} \emph{socorrer}, once more
  \item \emph{Say it:} a sentence of your own with \emph{socorrer}
  \item \emph{Recall:} say the Portuguese for the ambulance, then the Portuguese for the siren, then say \emph{socorrer} again
\end{itemize}

\begin{tcolorbox}[breakable,colback=teal!4,colframe=teal!35!black,title={Before you move on}]
What does socorrer mean? (\textquotedblleft{}To help (someone in danger), to come to the aid of\textquotedblright{}.) Say it once more.
\end{tcolorbox}
\section[o socorrista]{o socorrista — the first-aider, the paramedic}

\label{lesson:PT-C305-socorrista}



\begin{quote}
Before the new one: say the Portuguese for the siren, then the Portuguese for to help.
\end{quote}

\subsection*{You'll want to know: o socorrista}
\textbf{o socorrista} — \textquotedblleft{}the first-aider, the paramedic\textquotedblright{}.

\begin{quote}
O socorrista falou com calma ao rapaz.
\end{quote}

\begin{quote}
\emph{The paramedic spoke calmly to the young man.}
\end{quote}

The same form for a woman: \emph{a socorrista}.

\begin{cousinweb}
\emph{socorrista} is built on \emph{socorro} \textquotedblleft{}help\textquotedblright{} with \emph{-ista} for the person; \emph{socorro} goes back to Latin \emph{succurrere} \textquotedblleft{}to run to help\textquotedblright{} (EtymDB).
\end{cousinweb}

\begin{grammarlens}[title={What you've built}]
One more word for telling, explaining and describing at B1.
\end{grammarlens}

\subsection*{Your turn}
\begin{itemize}
\raggedright
  \item \emph{Say it:} \emph{o socorrista}
  \item \emph{Say it:} \emph{o socorrista}, once more
  \item \emph{Say it:} a sentence of your own with \emph{o socorrista}
  \item \emph{Recall:} say the Portuguese for the siren, then the Portuguese for to help, then say \emph{o socorrista} again
\end{itemize}

\begin{tcolorbox}[breakable,colback=teal!4,colframe=teal!35!black,title={Before you move on}]
What does o socorrista mean? (\textquotedblleft{}The first-aider, the paramedic\textquotedblright{}.) Say it once more.
\end{tcolorbox}
\section[a maca]{a maca — the stretcher}

\label{lesson:PT-C305-maca}



\begin{quote}
Before the new one: say the Portuguese for to help, then the Portuguese for the first-aider.
\end{quote}

\subsection*{You'll want to know: a maca}
\textbf{a maca} — \textquotedblleft{}the stretcher\textquotedblright{}.

\begin{quote}
Levaram o homem numa maca.
\end{quote}

\begin{quote}
\emph{They carried the man on a stretcher.}
\end{quote}

Not to be confused with \emph{a maçã}, the apple.

\begin{cousinweb}
\emph{maca} came from Spanish \emph{hamaca} \textquotedblleft{}a hammock\textquotedblright{} (EtymDB), a word taken from Taíno \emph{hamaka}.
\end{cousinweb}

\begin{grammarlens}[title={What you've built}]
One more word for telling, explaining and describing at B1.
\end{grammarlens}

\subsection*{Your turn}
\begin{itemize}
\raggedright
  \item \emph{Say it:} \emph{a maca}
  \item \emph{Say it:} \emph{a maca}, once more
  \item \emph{Say it:} a sentence of your own with \emph{a maca}
  \item \emph{Recall:} say the Portuguese for to help, then the Portuguese for the first-aider, then say \emph{a maca} again
\end{itemize}

\begin{tcolorbox}[breakable,colback=teal!4,colframe=teal!35!black,title={Before you move on}]
What does a maca mean? (\textquotedblleft{}The stretcher\textquotedblright{}.) Say it once more.
\end{tcolorbox}
